\documentclass[11pt,a4paper]{article}
\usepackage[margin=1in]{geometry}
\usepackage{amsmath,amssymb,booktabs}
\usepackage[hidelinks]{hyperref}
\raggedbottom
\widowpenalty=10000
\clubpenalty=10000
\setlength{\emergencystretch}{2em}
\newcommand{\Mpl}{M_{\rm Pl}}
\newcommand{\GN}{G_N}
\begin{document}
\title{Radial Dimensionality Theory VI: A Field-Theoretic Origin for the Density-Dependent Gravitational Coupling}
\author{Christopher K. Merrill\\\small Independent Researcher\\[1ex]
Claude (Anthropic)\\\small Large Language Model, Computational Collaboration\\[1ex]
Astra (Open AI)\\\small Large Language Model, Computational Collaboration}
\date{Consistency revision: November 12, 2026}
\maketitle
\begin{abstract}
We audit a scalar-coupling candidate motivated by the RDT density law.
For the written coupling $F(\sigma)=M_{\rm Pl}^2(1-\sigma)$, variation of
the action gives $M^2\Box\sigma=U'+M_{\rm Pl}^2R/2$. The previously chosen
increasing potential cannot support the claimed positive-density adiabatic
equilibrium on $0<\sigma<A$ with positive curvature. An exact constant-field
trace check gives the same obstruction for nonrelativistic matter and
nonnegative potential. This invalidates the proposed equilibrium derivation,
not the separately specified divergence-operator model. We retain the action
as an audited candidate, distinguish local coupling from the weighted radial
operator, and withdraw unsupported stellar, screening and completion claims.
\end{abstract}

\section*{Version note: v2}
This version retains the published title, but does not establish the field-theoretic origin named in it. The proposed positive-density equilibrium fails the equations of the written scalar action, so the revision audits a failed candidate derivation and withdraws the claimed origin and equivalence between local coupling and the weighted radial operator. The action remains as an audited candidate; earlier local-coupling stellar diagnostics are not solutions of the full scalar stellar problem. This failure does not exclude the separately specified phenomenological operator model or every possible field-theoretic completion.


\section{Three different prescriptions}
Let $x=(\rho/\rho_0)^\alpha$ and $\delta=-3Ax/(1+x)$. The operator used in
the series has radial weight
\begin{equation}
 W(r)=r^2(r/r_\star)^{\delta(r)},\qquad
 a_{\rm op}(r)=\frac{4\pi\GN}{W(r)}\int_0^r W(s)\rho(s)\,ds.
\end{equation}
The dimensionless ratio $r/r_\star$ is required when the exponent varies.
The local assignment $G_F=\GN/(1-\sigma)$ with $\sigma=Ax/(1+x)$ is a
separate prescription. Multiplying a TOV pressure gradient by $G_F/\GN$
is a third prescription: it is not a solution of the scalar equations.

Even for constant $\delta$, local equality fails for a regular nonuniform
profile. For $\rho(r)=\rho_c[1-q(r/R)^2]$,
\begin{equation}
 \frac{a_{\rm op}}{a_N}=
 \frac{1/(3+\delta)-q(r/R)^2/(5+\delta)}{1/3-q(r/R)^2/5}.
\end{equation}
At $q=0.7$, $r=R$, $\delta=-0.02$ this is $1.00866321$, whereas the
central ratio $3/(3+\delta)$ is $1.00671141$. This analytic counterexample
refutes universal local equivalence; it is not a fitted stellar model or a
no-go theorem for every possible covariant completion.

\section{Normalized action and field equations}
Use signature $(-+++)$, $\hbar=c=1$, $\Mpl^2=(8\pi\GN)^{-1}$ and
$T_{\mu\nu}=-2(\delta S_m/\delta g^{\mu\nu})/\sqrt{-g}$:
\begin{equation}
 S=\int\!d^4x\sqrt{-g}\left[\frac{F(\sigma)}2R
 -\frac{M^2}2(\nabla\sigma)^2-U(\sigma)\right]+S_m[g,\psi],
 \quad F=\Mpl^2(1-\sigma).
\end{equation}
Here $\sigma$ is dimensionless, $F$ and $M^2$ have mass dimension two,
and $U$ has mass dimension four. $M$ is a free normalization scale;
no screened value has been established. Direct variation gives
\begin{align}
 F G_{\mu\nu}&=T_{\mu\nu}+M^2\left(\nabla_\mu\sigma\nabla_\nu\sigma
 -\tfrac12g_{\mu\nu}(\nabla\sigma)^2\right)-g_{\mu\nu}U
 +\nabla_\mu\nabla_\nu F-g_{\mu\nu}\Box F,\\
 M^2\Box\sigma&=U'-\tfrac12F'R=U'+\tfrac12\Mpl^2R.
\end{align}
The metric derivative contribution is
$-\Mpl^2(\nabla_\mu\nabla_\nu\sigma-g_{\mu\nu}\Box\sigma)$;
the earlier schematic expression had the opposite sign. $G_F=1/(8\pi F)$
is the curvature coefficient's coupling, not automatically the gravitational
constant measured in an apparatus with scalar exchange.

\section{Why the proposed equilibrium fails}
The earlier construction assigned, for $\lambda>0$ of mass dimension four,
\begin{equation}
 U'=\lambda\left[\frac{\sigma}{A-\sigma}\right]^{1/\alpha},\qquad
 U(\sigma)=\lambda\int_0^\sigma
 \left[\frac{s}{A-s}\right]^{1/\alpha}ds.
\end{equation}
Solving the separately assigned equation $U'=\lambda\rho/\rho_0$ does
algebraically reproduce $\sigma=Ax/(1+x)$. However, with positive curvature
the action instead requires $U'=-\Mpl^2R/2$ for a constant field. Both
terms in $U'+\Mpl^2R/2$ are positive in the proposed matter branch.
The successful inversion in the historical \texttt{u\_sigma.py} tested an
assigned source equation with the wrong sign for this action.

This does not rely solely on substituting the GR trace. The exact trace is
\begin{equation}
 FR=-T+M^2(\nabla\sigma)^2+4U-3\Mpl^2\Box\sigma.
\end{equation}
At constant field with nonrelativistic $T=-\rho$, the scalar equation becomes
\begin{equation}
 0=U'+\frac{\rho+4U}{2(1-\sigma)}.
\end{equation}
For $\rho>0$, $0<\sigma<A<1$ and the specified $U\ge0$, there is no
solution. This excludes the claimed adiabatic branch, not arbitrary
pressure-supported, gradient-dominated or differently normalized theories.
Changing the coupling, potential or source is a model change requiring its
own analysis; no repair is silently substituted here.

Two further mathematical corrections matter. For $\alpha=1$ the stated
$U'=\lambda\sigma/(A-\sigma)$ is not globally linear: the potential is
quadratic only near zero. For $\alpha\le1$ the integral diverges at the
wall; a finite integral ending at $0.999A$ is not the value at $A$.
For $\alpha<1$, $U''(0^+)=0$, so the proposed bare potential alone does
not give a heavy vacuum scalar.

\section{What the kinetic and tensor checks establish}
Under $g^E_{\mu\nu}=(1-\sigma)g_{\mu\nu}$, the Einstein-frame scalar
kinetic coefficient (dimension mass squared) is
\begin{equation}
 K(\sigma)=\frac{M^2}{1-\sigma}
 +\frac{3\Mpl^2}{2(1-\sigma)^2},\qquad
 \frac{d\chi}{d\sigma}=\sqrt K.
\end{equation}
It is positive for $M^2>0$, $\sigma<1$. This is a kinetic-sign check,
not a proof of stable screened backgrounds. The canonical dimensionless
matter coupling is
\begin{equation}
 \beta_\chi=\Mpl\frac{d\ln(1-\sigma)^{-1/2}}{d\chi}
 =\frac{\Mpl}{2(1-\sigma)\sqrt K}.
\end{equation}
It depends on $M/\Mpl$; the previous value $1/2$ differentiated with respect
to the unnormalized field and was not the physical coupling. A mandatory
short-range branch does not follow from that value.

The two-derivative $F(\sigma)R$ action has tensor characteristics on the
metric null cone for a regular $F>0$ background \cite{Baker}; this says nothing about
scalar charge amplitudes or observed waveform agreement. A Ricci-flat
constant $\sigma=0$ background is possible if $U(0)=U'(0)=0$.
Additional scalar perturbations remain; a GR background does not establish
GR for all vacuum dynamics or all boundary conditions.

\section{Stellar and empirical status}
The canonical operator-model scan~\cite{Series_paper1} uses $\rho_0=150\,{\rm g\,cm^{-3}}$
and $\alpha\in[0.02,0.5]$. Its nominal solar proxy threshold, denoted
$A_{\rm lim}$, has a sampled and locally refined amplitude peak
$A=0.0087502019$ at $\alpha=0.1244971139$ (asymptote $3-3A=2.973749394$).
It has neither calibrated confidence coverage nor a certified continuous
maximum. Changing the physical response requires a new solar comparison;
sharing a density-law shape does not transfer this threshold.

The historical \texttt{geff\_check.py} used $(A,\alpha)=(0.0087,0.122)$,
not the manuscript's rounded $\alpha=0.12$, and the older \texttt{ns\_v7}
implementation. Its reported $-0.60\%$ mass and $-0.34\%$ radius shifts
are archived local-TOV diagnostics, not checked scalar-star predictions.
They are removed from the current prediction table. No new scalar stellar
integration is performed in this revision.

The current III--IV revisions~\cite{Series_paper3,Series_paper4} use CompOSE RG(SLy4), finite-pressure
cutoffs and separately checked phenomenological embeddings. Their modified
mass curves reach the EOS boundary without a bracketed maximum; their
conditional tidal extraction retains an approximately $1.30\%$
interior/exterior metric-gradient mismatch. None of these checks validates
this action or transfers a tidal shift to it. The old scalar-tensor companion~\cite{Series_companion}
uses a different action and is audited separately.

\section{Conclusion and version changes}
The local equilibrium derivation fails a necessary sign check. The operator
series remains a separate empirical model, and the local Poisson ansatz can
also be tested as a declared phenomenological model. Neither is derived by
the action as currently specified. Claims of completion, screening,
automatic solar-bound transfer and reproduced tides are withdrawn.
The next theoretical task is an explicit revised action/source construction,
followed by its own field and observational calculations.
\section*{Code and revision availability}
Producing scripts, inputs, saved results and numerical audit evidence are
publicly available in the \href{https://github.com/chrismerrill1974/solarcore_code/tree/main/reproduction}{shared RDT reproducibility archive}.
The accompanying source package, \nolinkurl{merrill2026_field_theoretic_rdt_origin_code}, includes the
manuscript, dependencies and resource-controlled reproduction commands.
Its \texttt{reproduction/CAPSULE.json} records file hashes;
\texttt{reproduce.py} verifies the package and rebuilds the paper from retained
scientific outputs. Full scientific runs are separate from this saved-result
reproduction. Historical scripts and results remain separately identified.

\paragraph{Computational collaboration.} Claude (Anthropic) and Astra (Open AI).

\begin{thebibliography}{9}
\bibitem{Baker} T.~Baker et al., ``Strong constraints on cosmological gravity
from GW170817 and GRB 170817A,'' Phys. Rev. Lett. \textbf{119}, 251301 (2017),
\href{https://arxiv.org/abs/1710.06394}{arXiv:1710.06394}.
\bibitem{Series_companion} C.~K.~Merrill, Claude (Anthropic), and Astra (Open AI),
``A Scalar-Tensor Companion to Radial Dimensionality Theory: A Chameleon Model Reproducing the Density Law, with Gravitational-Wave Predictions,''
revised version, November 12, 2026.

\bibitem{Series_paper1} C.~K.~Merrill, Claude (Anthropic), and Astra (Open AI),
``Radial Dimensionality Theory I: Solar Neutrino and Helioseismic Constraints on Density-Dependent Spatial Dimensionality,''
revised version, November 12, 2026.

\bibitem{Series_paper3} C.~K.~Merrill, Claude (Anthropic), and Astra (Open AI),
``Radial Dimensionality Theory III: Neutron Star Structure under Solar-Constrained Density-Dependent Dimensionality,''
revised version, November 12, 2026.

\bibitem{Series_paper4} C.~K.~Merrill, Claude (Anthropic), and Astra (Open AI),
``Radial Dimensionality Theory IV: Tidal Deformability Ceilings under Solar-Constrained Density-Dependent Dimensionality,''
revised version, November 12, 2026.

\end{thebibliography}
\end{document}
